\section{Discussion and limitations}\label{sec:discussion}
The three substantive questions have distinct answers in this bounded
study. Independent natural-source confirmation does not reproduce the
large development advantage over raw RBF. Complete measured costs make
the confirmed accuracy--cost comparison less favorable to raw-plus-KAHM,
while simpler raw ridge remains a strong alternative. Finally, the
specified geometric diagnostic increment supplies no improvement in
confirmation decisions beyond metadata and a paid validation probe.

These results are useful because the tests could have supported a
different conclusion. The frozen design retains conditional improvements,
including the small mean gain over raw RBF on Saugeen, as well as harms
and successful controls. It does not infer a universal failure from four
reserved families. The contribution is a reproducible examination of
where evidence for this particular KAHM forecasting procedure weakens:
a favorable development aggregate is concentrated, structural controls
explain much of the apparent benefit, and the proposed descriptors do
not yield a transferable adoption rule. A stronger claim about KAHM
would require further independent evidence, with a new confirmation
collection reserved before any new diagnostic is developed.


\subsection{What a representation comparison establishes}
Raw augmentation gives a clear information baseline: the original history
remains recoverable from the supplied code. A forecasting gain therefore
concerns the trained finite decoder and its regularization. It does not show
that KAHM has created additional target information. Conversely, a loss from
augmentation does not refute information preservation. Additional standardized
coordinates change distances, the finite Nystr\"om basis, and the effective
regularization geometry. These changes can make the prediction procedure
better or worse even though the raw coordinates remain present.

The affine and locality-aware augmentations are essential controls for this
reason. A favorable raw-versus-KAHM comparison by itself cannot attribute the
benefit to nonlinear KAHM reconstruction. Local coordinates, repeated views
of raw inputs, and their interaction with the downstream feature basis can
also matter. The declared alternatives share reference-selection rules and
have comparable local-coordinate constructions, while being fitted as
independent computational procedures. Actual feature widths and retained
Nystr\"om ranks are reported instead of assuming that the same landmark
count implies identical effective capacity.

\subsection{Limits of predicting benefit}
The diagnostic experiment concerns four specified, low-dimensional ridge
rules and one fixed reduced-capacity probe. A failure of those rules does
not establish that usefulness is unpredictable. A successful decision on
one family would likewise not establish a general adoption criterion.
With four development families, a complex learned selector would be hard
to validate independently; the present rules deliberately restrict their
flexibility and retain every prespecified variant.

Two design limits are especially relevant. The descriptors do not directly
encode the downstream landmark count. When the training-window cap does
not bind, the two capacity settings have identical diagnostic inputs even
though their full predictors differ. The history and horizon are also
changed together, so their individual effects cannot be identified by this
experiment. These are limitations of the evaluated criterion, and should
inform a future independently evaluated diagnostic rather than a retuning
of the reserved-source results.

Furthermore, predicting conditional mean log benefit is not the same
decision problem as minimizing conditional expected scaled MSE. Although
the sign of each realized log ratio agrees with the sign of its loss
difference, expectations of these quantities need not have the same sign
under heterogeneous losses. We therefore report both estimated-benefit
errors and the actual loss and cost of the resulting decisions. Input
geometry alone also cannot determine a universally preferable predictor:
different target relationships can be associated with the same input law.
The paid validation probe supplies target-relevant evidence, and the
geometric increment is evaluated against that information explicitly.

\subsection{Data, statistical scope, and computational scope}
The eight source families were deliberately selected from a public archive
for diversity and feasibility. They are not a probability sample of future
applications. Four families per phase support bounded descriptive comparisons;
series, horizons, capacity settings, seeds, and overlapping windows do not
increase that source count. We report family effects and sensitivity to
leaving out one family without attaching population confidence claims to
the small collection.

Some records are short. The longer Hospital and Traffic settings have only
21 and 29 training windows, respectively. NN5 T88 has seven available test
origins in the longer setting. The fixed roster retains this case, and its
loss estimates should be read with that limited support in mind. Missing
windows are excluded using the same availability rule for every method;
the resulting estimand concerns available complete windows, rather than
unobserved targets or operational forecasts with missing contexts. The
FRED-MD experiment is relative to a supplied retrospective snapshot and
cannot certify every upstream transformation or emulate real-time release
and revision constraints.

The comparison is restricted to local univariate forecasters, two small
capacity settings, and finite tuning grids. It includes a small MLP and
boosted-tree control but no global pretrained forecaster or comprehensive
state-of-the-art neural benchmark. The training-window cap is nonbinding
for the development sources and binding for the long confirmation series;
the larger setting therefore changes available sample count only for part
of the collection. The three-value alpha grid is fixed in advance, and a
boundary choice is not evidence that an optimum over all regularization
strengths has been located.

Measured costs apply to the supplied reference implementations on the
recorded CPU environment. Warm-library scientific-pipeline costs include
required feature construction, selection, refitting, prediction, and
serialization; library initialization and experimental process launch are
separate measurements. Timing small procedures can be sensitive to system
load and Python bookkeeping. The component records expose those effects.
Process high-water memory and serialized model size answer different
questions. The former includes the interpreter and input arrays and is
not an incremental memory estimate for a hypothetical combined selector.

\section{Conclusion}\label{sec:conclusion}
We evaluated KAHM-derived local features with independent natural-source
confirmation, explicit computational dependencies, and a prospectively
frozen usefulness diagnostic. The preserved synthetic study motivated
the experiment but was not counted as natural confirmation. On the four
reserved families, raw-plus-KAHM has higher aggregate error and cost
than the matched raw RBF predictor, and raw ridge has lower aggregate error than
either. Advantages over affine/locality augmentation are small and
source-dependent. Adding geometric features to metadata and a validation
probe changes no confirmation choices and provides no reduction in
log-benefit prediction error. The supplied artifacts preserve the
failures as well as improvements and independently reconstruct every
reported task--method loss. The evidence supports evaluating KAHM as
one conditional representation choice within a complete forecasting
procedure; it does not establish a broadly preferable predictor or a
validated rule for anticipating its benefit.

